\section{A direct state-map formulation}
\label{sec:direct-state-map}

The return fold also admits a stepwise formulation with an exact time
change. Define the state map $f$ by $f(27)=15$ and $f(q)=q$ otherwise,
and extend it to configurations without changing the head or tape.
Strengthen the old-machine invariant to
\[
 J(q,h,\tau)\ \Longleftrightarrow\
 \bigl(q\in\{15,18,47,65\}\Rightarrow\tau(h-1)=0\bigr)
 \ \land\
 \bigl(q=27\Rightarrow\tau(h)=0\bigr).
\]
The first conjunct was established by the incoming write-zero/right
inventory. The only entries to $27$ are the read-zero left moves from
$18$, $47$, and $65$; the first conjunct therefore establishes the
second at every such entry. Both conjuncts are initially vacuous at
state $95$. Thus $J$ is inductive for the old machine.

For every old transition row other than $(15,0)$ and $(27,1)$, the exact
finite table identity is
\[
 \delta_N(f(q),b)=(w,d,f(r))
 \qquad\text{whenever}\qquad
 \delta_O(q,b)=(w,d,r).
\]
There are $484$ such row identities. The row $(27,1)$ is infeasible by
$J$. The remaining row $(15,0)$ is implemented by the three-step
return block through $32$ and $45$, using the two adjacent zero cells
provided by $J$ and the read-zero guard. Therefore one old step from any
reachable configuration has exactly the same mapped full endpoint
after one new step, except that old $(15,0)$ takes three new steps.
Neither of the two intermediate states in this longer block is halting.

For an old run prefix of length $n$, let $k(n)$ count its transitions
using row $(15,0)$, and set
\[
 \theta(n)=n+2k(n).
\]
Induction on $n$ proves
\[
 \operatorname{Run}_N\bigl(\theta(n),f(c_0)\bigr)
 =f\bigl(\operatorname{Run}_O(n,c_0)\bigr)
\]
whenever the old run has not halted before step $n$. Because
$\theta(n+1)-\theta(n)$ is either one or three, this aligns unbounded
runs as well as finite prefixes. The state map preserves and reflects
halt, and the inserted states are nonhalting. Thus either both machines
halt or neither does. If the old machine halts after exactly $T$ steps,
then the new machine halts after exactly
\[
 T+2k(T),\qquad\text{in particular}\qquad
 T\le T_{\rm new}\le3T.
\]
This exact time-change theorem holds for every initial tape and head
at the shared start state $95$. Its proof is independent of the logical
consistency calibration. It is an alternative proof of the same binary
machine comparison, not an additional state reduction.
